\documentclass[12pt]{article}

% Idioma y codificación
\usepackage[spanish, es-tabla]{babel}       %es-tabla para que se titule "Tabla"
\usepackage[utf8]{inputenc}

% Márgenes
\usepackage[a4paper,top=3cm,bottom=2.5cm,left=3cm,right=3cm]{geometry}

% Comentarios de bloque
\usepackage{verbatim}

% Paquetes de links
\usepackage[hidelinks]{hyperref}    % Permite enlaces
\usepackage{url}                    % redirecciona a la web

% Más opciones para enumeraciones
\usepackage{enumitem}


%------------------------------------------------------------------------

%Paquetes de figuras
\usepackage{caption}
\usepackage{subcaption} % Figuras al lado de otras
\usepackage{float}      % Poner figuras en el sitio indicado H.


% Paquetes de imágenes
\usepackage{graphicx}       % Paquete para añadir imágenes
\usepackage{transparent}    % Para manejar la opacidad de las figuras

% Paquete para usar colores
\usepackage[dvipsnames]{xcolor}
\usepackage{pagecolor}      % Para cambiar el color de la página

% Habilita tamaños de fuente mayores
\usepackage{fix-cm}


%------------------------------------------------------------------------

% Paquetes de matemáticas
\usepackage{amsmath, amsfonts, amssymb, mathrsfs}
\usepackage{cancel}     % Simplificar tachando
\usepackage{polynom}    % Divisiones y Ruffini

\usepackage{pgfplots}   %Representar funciones
\pgfplotsset{compat=1.18}  % Versión 1.18

%Definición de teoremas, etc.
\usepackage{amsthm}
%\swapnumbers   % Intercambia la posición del texto y de la numeración

\theoremstyle{plain}
\newtheorem{teo}{Teorema}[chapter]     % Se resetea en cada chapter
\newtheorem{coro}{Corolario}[teo]           % Se resetea en cada teorema
\newtheorem{prop}[teo]{Proposición}         % Usa el mismo contador que teorema
\newtheorem{lema}[teo]{Lema}                % Usa el mismo contador que teorema

\theoremstyle{remark}
\newtheorem*{observacion}{Observación}

\theoremstyle{definition}
\newtheorem{definicion}{Definición} [chapter]     % Se resetea en cada chapter
\newtheorem*{notacion}{Notacion}
\newtheorem*{ejemplo}{Ejemplo}
\newtheorem*{ejercicio*}{Ejercicio}             % No numerado
\newtheorem{ejercicio}{Ejercicio} [section]     % Se resetea en cada section


% Modificar el formato de la numeración del teorema "ejercicio"
\renewcommand{\theejercicio}{%
  \ifnum\value{section}=0 % Si no se ha iniciado ninguna sección
    \arabic{ejercicio}% Solo mostrar el número de ejercicio
  \else
    \thesection.\arabic{ejercicio}% Mostrar número de sección y número de ejercicio
  \fi
}


% \renewcommand\qedsymbol{$\blacksquare$}         % Cambiar símbolo QED
%------------------------------------------------------------------------

% Paquetes para encabezados
\usepackage{fancyhdr}
\pagestyle{fancy}
\fancyhf{}

\newcommand{\helv}{ % Modificación tamaño de letra
\fontfamily{}\fontsize{12}{12}\selectfont}
\setlength{\headheight}{15pt} % Amplía el tamaño del índice


%\usepackage{lastpage}   % Referenciar última pag   \pageref{LastPage}
\fancyfoot[C]{\thepage}

%------------------------------------------------------------------------

\begin{comment} % Descomentar en book. Comentar en article

% Conseguir que no ponga "Capítulo 1". Sino solo "1."
\makeatletter
\renewcommand{\@makechapterhead}[1]{%
  \vspace*{50\p@}%
  {\parindent \z@ \raggedright \normalfont
    \ifnum \c@secnumdepth >\m@ne
      \huge\bfseries \thechapter.\hspace{1em}\ignorespaces
    \fi
    \interlinepenalty\@M
    \Huge \bfseries #1\par\nobreak
    \vskip 40\p@
  }}
\makeatother
\end{comment}

%------------------------------------------------------------------------
% Personalizar la portada
\usepackage{titling}


%------------------------------------------------------------------------

% Comandos definidos
\newcommand{\bb}[1]{\mathbb{#1}}
\newcommand{\cc}[1]{\mathcal{#1}}
\newcommand{\Hop}{\stackrel{L'H\hat{o}pital}{=}}

\newcommand{\red}[1]{{\color{red}#1}} % Para integrales, destacar los cambios.

% Método de integración
\newcommand{\MetInt}[2]{
    \left[\begin{array}{c}
        #1 \\ #2
    \end{array}\right]
}

%------------------------------------------------------------------------


\begin{document}

    % 1. Foto de fondo
    % 2. Título
    % 3. Encabezado Izquierdo
    % 4. Color de fondo
    % 5. Coord x del titulo
    % 6. Coord y del titulo
    % 7. Fecha

    
    % 1. Foto de fondo
% 2. Título
% 3. Encabezado Izquierdo
% 4. Color de fondo
% 5. Coord x del titulo
% 6. Coord y del titulo
% 7. Fecha

\newcommand{\portada}[7]{

    \portadaBase{#1}{#2}{#3}{#4}{#5}{#6}{#7}
    \portadaBook{#1}{#2}{#3}{#4}{#5}{#6}{#7}
}

\newcommand{\portadaExamen}[7]{

    \portadaBase{#1}{#2}{#3}{#4}{#5}{#6}{#7}
    \portadaArticle{#1}{#2}{#3}{#4}{#5}{#6}{#7}
}




\newcommand{\portadaBase}[7]{

    % Tiene la portada principal y la licencia Creative Commons
    
    % 1. Foto de fondo
    % 2. Título
    % 3. Encabezado Izquierdo
    % 4. Color de fondo
    % 5. Coord x del titulo
    % 6. Coord y del titulo
    % 7. Fecha
    
    
    \thispagestyle{empty}               % Sin encabezado ni pie de página
    \newgeometry{margin=0cm}        % Márgenes nulos para la primera página
    
    
    % Encabezado
    \fancyhead[L]{\helv #3}
    \fancyhead[R]{\helv \nouppercase{\leftmark}}
    
    
    \pagecolor{#4}        % Color de fondo para la portada
    
    \begin{figure}[p]
        \centering
        \transparent{0.3}           % Opacidad del 30% para la imagen
        
        \includegraphics[width=\paperwidth, keepaspectratio]{assets/#1}
    
        \begin{tikzpicture}[remember picture, overlay]
            \node[anchor=north west, text=white, opacity=1, font=\fontsize{60}{90}\selectfont\bfseries\sffamily, align=left] at (#5, #6) {#2};
            
            \node[anchor=south east, text=white, opacity=1, font=\fontsize{12}{18}\selectfont\sffamily, align=right] at (9.7, 3) {\textbf{\href{https://losdeldgiim.github.io/}{Los Del DGIIM}}};
            
            \node[anchor=south east, text=white, opacity=1, font=\fontsize{12}{15}\selectfont\sffamily, align=right] at (9.7, 1.8) {Doble Grado en Ingeniería Informática y Matemáticas\\Universidad de Granada};
        \end{tikzpicture}
    \end{figure}
    
    
    \restoregeometry        % Restaurar márgenes normales para las páginas subsiguientes
    \pagecolor{white}       % Restaurar el color de página
    
    
    \newpage
    \thispagestyle{empty}               % Sin encabezado ni pie de página
    \begin{tikzpicture}[remember picture, overlay]
        \node[anchor=south west, inner sep=3cm] at (current page.south west) {
            \begin{minipage}{0.5\paperwidth}
                \href{https://creativecommons.org/licenses/by-nc-nd/4.0/}{
                    \includegraphics[height=2cm]{assets/Licencia.png}
                }\vspace{1cm}\\
                Esta obra está bajo una
                \href{https://creativecommons.org/licenses/by-nc-nd/4.0/}{
                    Licencia Creative Commons Atribución-NoComercial-SinDerivadas 4.0 Internacional (CC BY-NC-ND 4.0).
                }\\
    
                Eres libre de compartir y redistribuir el contenido de esta obra en cualquier medio o formato, siempre y cuando des el crédito adecuado a los autores originales y no persigas fines comerciales. 
            \end{minipage}
        };
    \end{tikzpicture}
    
    
    
    % 1. Foto de fondo
    % 2. Título
    % 3. Encabezado Izquierdo
    % 4. Color de fondo
    % 5. Coord x del titulo
    % 6. Coord y del titulo
    % 7. Fecha


}


\newcommand{\portadaBook}[7]{

    % 1. Foto de fondo
    % 2. Título
    % 3. Encabezado Izquierdo
    % 4. Color de fondo
    % 5. Coord x del titulo
    % 6. Coord y del titulo
    % 7. Fecha

    % Personaliza el formato del título
    \pretitle{\begin{center}\bfseries\fontsize{42}{56}\selectfont}
    \posttitle{\par\end{center}\vspace{2em}}
    
    % Personaliza el formato del autor
    \preauthor{\begin{center}\Large}
    \postauthor{\par\end{center}\vfill}
    
    % Personaliza el formato de la fecha
    \predate{\begin{center}\huge}
    \postdate{\par\end{center}\vspace{2em}}
    
    \title{#2}
    \author{\href{https://losdeldgiim.github.io/}{Los Del DGIIM}}
    \date{Granada, #7}
    \maketitle
    
    \tableofcontents
}




\newcommand{\portadaArticle}[7]{

    % 1. Foto de fondo
    % 2. Título
    % 3. Encabezado Izquierdo
    % 4. Color de fondo
    % 5. Coord x del titulo
    % 6. Coord y del titulo
    % 7. Fecha

    % Personaliza el formato del título
    \pretitle{\begin{center}\bfseries\fontsize{42}{56}\selectfont}
    \posttitle{\par\end{center}\vspace{2em}}
    
    % Personaliza el formato del autor
    \preauthor{\begin{center}\Large}
    \postauthor{\par\end{center}\vspace{3em}}
    
    % Personaliza el formato de la fecha
    \predate{\begin{center}\huge}
    \postdate{\par\end{center}\vspace{5em}}
    
    \title{#2}
    \author{\href{https://losdeldgiim.github.io/}{Los Del DGIIM}}
    \date{Granada, #7}
    \thispagestyle{empty}               % Sin encabezado ni pie de página
    \maketitle
    \vfill
}
    \portadaExamen{ffccA4.jpg}{Geometría III\\Examen IV}{Geometría III. Examen IV}{MidnightBlue}{-8}{28}{2023-2024}{Arturo Olivares Martos}

    \begin{description}
        \item[Asignatura] Geometría III.
        \item[Curso Académico] 2023-24.
        \item[Grado] Grado en Matemáticas.
        \item[Grupo] B.
        \item[Profesor] José María Espinar García.
        \item[Descripción] Parcial de los Temas 1 y 2.
        \item[Fecha] 7 de noviembre de 2023.
        %\item[Duración] 90 minutos.
    
    \end{description}
    \newpage
    
    \begin{ejercicio}
        Sean $R_1$ y $R_2$, $R_1\cap R_2=\emptyset$, dos rectas paralelas en un plano afín $\cc{A}$.
        Sean $a_i,b_i,c_i\in R_i$, $i=1,2$, tres puntos distintos en cada una de las rectas.
        Demostrar que si $R_{a_1b_2}\|R_{a_2b_1}$ y $R_{b_1c_2}\|R_{b_2c_1}$, entonces $R_{a_1c_2}\|R_{a_2c_1}$.\\

        Es la demostración del Teorema de Pappus en el caso de que las rectas sean paralelas.
    \end{ejercicio}

    \begin{ejercicio}
        Sea $T:=\{a_1,a_2,a_3\}$ un triángulo en un plano afín $\cc{A}$. Denotemos por $a_i'$, $i=1,2,3$, al punto medio del dado opuesto a $a_i$, es decir, $a_i'=a_j+\frac{1}{2}\vec{a_ja_k}$, $i\neq j\neq k$ e $i,j,k\in \{1,2,3\}$. Demostrar que existe una homotecia $h:\cc{A}\to \cc{A}$ tal que $h(a_i)=a_i'$ para $i=1,2,3$. Calcular el centro y la razón de dicha homotecia.\\

        En la demostración del Teorema de Euler se vio que dicha homotecia existe y que su centro es el baricentro del triángulo y su razón es $-\nicefrac{1}{2}$.
    \end{ejercicio}

    \begin{ejercicio}
        Sean $S_1:=\{(x,y,z)\in \bb{R}^3\mid x+2y+z=-1\}$ y $S_2=(1,2,0)+\cc{L}\{(-1,1,1), (0,0,1)\}$.
        \begin{enumerate}
            \item Calcula la intersección $I:=S_1\cap S_2$ y la suma.
            
            Calcularemos en primer lugar su intersección. Para ello, calculamos la ecuación implícita de $S_2$:
            \begin{equation*}
                S_2\equiv \begin{vmatrix}
                    -1 & 0 & x-1 \\
                    1 & 0 & y-2 \\
                    1 & 1 & z
                \end{vmatrix} = 0 = -\begin{vmatrix}
                    -1 & x-1 \\
                    1 & y-2
                \end{vmatrix} = y-2 + x-1 = y+x-3
            \end{equation*}

            Por tanto, calculamos ahora $I=S_1\cap S_2$:
            \begin{equation*}
                \left\{
                    \begin{array}{c}
                        x+2y+z=-1 \\
                        y+x-3=0
                    \end{array}
                \right\} \Longrightarrow \left\{
                    \begin{array}{l}
                        x=\lm +7\\
                        y=-\lm -4\\
                        z=\lm
                    \end{array}
                \right. \qquad \lm\in \bb{R}
            \end{equation*}

            Por tanto, $I=(7, -4, 0)+\cc{L}\{(-1, -1, 1)\}$. Como $S_1\cap S_2\neq \emptyset$, por la fórmula de las dimensiones, tenemos que:
            \begin{equation*}
                \dim S_1\vee \dim S_2 = \dim S_1 + \dim S_2 - \dim S_1\cap S_2 = 2+2-1=3
            \end{equation*}
            Además, como $S_1\vee S_2\subset \bb{R}^3$, tenemos que $S_1\vee S_2 = \bb{R}^3$.
            
            \item Sea $\cc{R}:=\{(1,-1,0), (0,0,1), (1,0,-1), (2,-2,1)\}$. Demostrar que es un sistema de referencia en $\bb{R}^3$.
            
            Para ello, mostraremos que su base asociada es base de $\bb{R}^3$. Esta es:
            \begin{equation*}
                \cc{B} = \left\{(-1, 1, 1), (0, 1, -1), (1, -1, 1)\right\}
            \end{equation*}

            Tenemos que:
            \begin{equation*}
                \begin{vmatrix}
                    -1 & 0 & 1 \\
                    1 & 1 & -1 \\
                    1 & -1 & 1
                \end{vmatrix} = -2 \neq 0
            \end{equation*}

            Por tanto, $\cc{B}$ es base de $\bb{R}^3$ y, por tanto, $\cc{R}$ es un sistema de referencia en $\bb{R}^3$.

            \item Sea el plano afín $P:= \{(\alpha,\beta,\gamma)_{\cc{R}}\in \bb{R}^3\mid \gamma=1\}$ dado en coordenadas respecto de $\cc{R}$. Calcular, si es posible, el ángulo de intersección entre $P$ e $I$. En caso de no ser posible, decir por qué.
            
            Tenemos que $P$ es un hiperplano e $I$ es una recta. Por tanto, sí se puede calcular el ángulo de intersección entre ambos.
            Para ello, necesitamos calcular en primer lugar un vector normal a $P$. Para ello, calculamos la ecuación implícita de $P$ en el sistema de referencia usual.
            Como $\gamma = 0+x-y-z$, tenemos que $P\equiv x-y-z=1$,
            por lo que un vector normal a $P$ es $v_P=(1,-1,-1)$.
            Por tanto, el ángulo de intersección entre $P$ e $I$ es:
            \begin{equation*}
                \theta  =\frac{\pi}{2}- \arccos\left(\frac{\langle (1,-1,-1)~(-1,-1,1)\rangle}{\sqrt{3}^2}\right)
                =\frac{\pi}{2}- \arccos\left(\frac{-1}{3}\right)
                \approx -0.34 \text{ rad}
            \end{equation*}
        \end{enumerate}
    \end{ejercicio}

    \begin{ejercicio}
        Sea $f:\bb{R}^3\to \bb{R}^3$ la aplicación afín dada por:
        \begin{gather*}
            f(-1,1,0) = (\nicefrac{11}{5}, \nicefrac{-2}{5}, 3) \qquad
            f(0,0,1) = (\nicefrac{18}{5}, \nicefrac{-1}{5}, 4) \\
            f(1,0,-1) = (\nicefrac{8}{5}, \nicefrac{-6}{5}, 4) \qquad
            f(2,-2,1)=(2,1,2)
        \end{gather*}
        \begin{enumerate}
            \item ¿Es $f$ una isometría de $\bb{R}^3$?
            \item Calcular el conjunto de puntos fijos de $f$ en el sistema de referencia usual.
            \item Calcula las ecuaciones que representan a $f$ respecto de los sistemas de referencia $\cc{R}$ (en el dominio) y $\cc{R}'$ (en el codominio); siendo:
            \begin{equation*}
                \begin{split}
                    \cc{R}&= \{(1,-1,0), (0,0,1), (1,0,-1), (2,-2,1)\} \\
                    \cc{R}'&= \{(0,0,1), (1,-1,0), (0,0,0), (-2,0,1)\} \\
                \end{split}
            \end{equation*}
        \end{enumerate}
    \end{ejercicio}   
\end{document}